\documentclass[12pt,a4paper]{article}
\usepackage[top=2.4cm,bottom=2.4cm,left=2.6cm,right=2.6cm]{geometry}
\usepackage{xeCJK}
\setCJKmainfont{Microsoft YaHei}
\setCJKsansfont{Microsoft YaHei}
\setmainfont{Times New Roman}
\usepackage{graphicx}
\usepackage{tikz}
\usetikzlibrary{positioning,arrows.meta,fit,backgrounds,shapes.geometric,chains,calc}
\usepackage{enumitem}
\usepackage{tabularx}
\usepackage{array}
\usepackage{booktabs}
\usepackage{hyperref}
\hypersetup{colorlinks=true,linkcolor=black,urlcolor=black}
\setlength{\parindent}{2em}

\title{\textbf{专利申请技术交底}\\[0.3em]
\large 人机共享记忆的分层智能索引与上下文按需下探方法}
\author{}
\date{\today}

\begin{document}
\maketitle
\thispagestyle{empty}

\noindent\textbf{前置说明}：本文为专利代理撰写正式申请文件的素材，按交底书习惯组织为技术领域、背景技术、要解决的技术问题、技术方案、实施方式、技术效果与创新点。文中量化数据均来自本产品真实实现与测试。

\section{技术领域}
本发明涉及人工智能与人机协同领域，具体涉及一种人机共享的项目记忆管理系统。人与 AI 助手在同一份图结构记忆上协作，系统用分层索引机制控制 AI 推理时的上下文成本，并保证双方写入的一致性。可用于 AI 编程助手的长程任务记忆、研发过程知识管理、多智能体共享工作记忆等场景。

\section{背景技术（现有技术及其缺陷）}
现有「AI 助手 + 项目记忆」类方案（向量数据库记忆、对话记录检索、RAG 上下文窗口等）普遍存在六类问题。

\begin{enumerate}[leftmargin=2em,itemsep=2pt]
\item \textbf{信息载体覆盖不全。} 这是 RAG 的普遍局限。现有 RAG / 向量库方案几乎全部以文本为中心，嵌入、召回、注入上下文的都是文本片段。图片、表格、图表、设计稿、专业文档（Word/PDF/Excel）等非文本信息要么被丢弃，要么被强行 OCR 或转文本，造成信息保真度损失。真实人机协作场景中，关键证据常常以原始文件格式存在。图片就是图片，Word 就是 Word，脱离原始载体即丢失信息。这些原始文件又往往由人手工协作产生，是 RAG 体系根本覆盖不到的盲区。
\item \textbf{上下文成本不可控。} 多数方案把记忆内容（对话历史、文档全文、向量片段）成批注入模型上下文。记忆规模增长时，单轮推理的上下文占用线性膨胀，既超窗口上限又稀释重点信息，token 成本失控。
\item \textbf{检索即全量扫描。} 无索引或索引与正文未分离时，每次查询都要把全部记忆文件读入内存做相似度或关键词匹配，机器 IO 与内存随记忆量增长，规模不可扩展。
\item \textbf{检索不可解释。} 向量检索（Embedding）的召回结果说不清为什么是这条，用户难以验证与纠偏；引入外部模型服务又违背本地优先、零依赖部署约束。
\item \textbf{人类输入丢失。} 人与 AI 共享同一份记忆时，人类通过正文或文档写入的关键信息没有未读信号，AI 只能靠关键词碰巧命中，信息容易被忽略；AI 也缺少已读确认机制，无法区分看过了和没看过。
\item \textbf{双写不一致。} 人与 AI 同时修改同一记忆对象时，静默覆盖会造成数据丢失；没有原子提交与冲突保护的方案存在崩溃丢数据风险。
\end{enumerate}

\section{要解决的技术问题}
提供一种分层记忆索引与按需下探的方法，使 AI 助手在共享记忆系统中的单轮推理上下文成本与记忆总量解耦，达到近乎常数级。同时满足五项要求：
\begin{enumerate}[leftmargin=2em,itemsep=2pt]
\item 索引与真实内容永远一致，人机双写环境下也不脱节；
\item 人类在任意位置的新输入不被 AI 遗漏；
\item 所有召回可解释；
\item 所有写入原子且幂等；
\item 任意格式的记忆对象以原始载体完整保真地可达，不因限于文本而丢失图片、文档等关键证据。
\end{enumerate}

\section{技术方案（核心技术思想）}
\subsection{总体架构：四级分层记忆}
将一份记忆按粒度与访问频率分成四层，层间以路径引用关联，如图~\ref{fig:layers} 所示。层级访问规则可以归纳为一句话，给 AI 的上下文只包含第 1 层全量、第 2 层命中摘要和路径，第 3、4 层进入上下文必须以工具调用显式下探为前提。

\begin{figure}[ht]
\centering
\begin{tikzpicture}[font=\small,every node/.style={align=left,inner sep=6pt},
  layer/.style={draw=black!70,fill=#1,rounded corners=2pt,minimum height=22pt}]
  \node[layer=blue!12] (l1) at (0,0) {\textbf{第 1 层 结构化索引（图形脑图）}\\节点/路线/方案线/标注的结构、状态与布局，KB 级，随上下文全量携带};
  \node[layer=green!12, below=4pt of l1] (l2) {\textbf{第 2 层 卡片索引（摘要+指纹）}\\标题、前 N 字摘要、指纹（mtime/size/etag）、附件清单；不存全文，检索只在这一层};
  \node[layer=yellow!22, below=4pt of l2] (l3) {\textbf{第 3 层 详情内容（Markdown 正文与补充文件）}\\按对象目录存放，仅在命中时被下探读取};
  \node[layer=orange!20, below=4pt of l3] (l4) {\textbf{第 4 层 资产文件（任意格式，对象资料包）}\\每个对象对应一个资料包目录，容纳任意格式附件；以路径+元数据契约暴露，读取能力由消费方自备};
  \node[right=22pt of l4, draw=black!60, dashed, rounded corners, minimum height=40pt, minimum width=95pt] (rule) {\textbf{访问规则}\\全量携带\\摘要+路径命中\\显式下探\\显式下探(路径)};
  \draw[->,thick] (l1.east) -- +(.9,0) |- (rule.west);
  \draw[->,thick] (l2.east) -- +(.9,0) |- (rule.west);
  \draw[->,thick] (l3.east) -- +(.9,0) |- (rule.west);
  \draw[->,thick] (l4.east) -- +(.9,0) |- (rule.west);
  \node[below=2pt of rule, font=\scriptsize, text=black!60] {路径引用关联};
\end{tikzpicture}
\caption{四级分层记忆架构与访问规则}\label{fig:layers}
\end{figure}

\subsection{卡片索引的生成与维护（人机双写一致性）}
\begin{itemize}[leftmargin=2em,itemsep=2pt]
\item \textbf{生成。} 任何正文文件的构建路径（创建对象、保存文档）完成后，同步为对象生成或刷新卡片，提取首标题与前 N 字摘要、计算内容哈希与文件指纹、枚举对象目录附件清单。索引以单个 JSON 文件原子写盘。
\item \textbf{校验用廉价指纹。} 查询时对每张卡只做一次 lstat，比较 mtime 和 size。指纹一致则直接用卡片摘要，实现零全文读取；不一致则只重读该文件并刷新单卡。
\item \textbf{双写刷新覆盖全部写路径。} 人类经应用保存、AI 经工具写入、文件改名、归档恢复、附件增删、对象创建，所有写路径统一在提交成功后触发单对象卡片刷新。
\item \textbf{绕桥修改自愈。} 用户用外部编辑器直接改文件（绕过系统）时，指纹失效机制在下次查询自动发现并只重读该文件，索引不会长期失真，如图~\ref{fig:heal} 所示。
\item \textbf{降级安全。} 索引文件损坏或缺失时自动按需重建，不会被陈旧索引阻塞读取。
\item \textbf{实测统计}（500 对象规模压测）：第二次查询的全文读取次数为 0；外部改动 1 个文件后的查询全文读取次数为 1；首次建索引为一次性全量扫描，次数约等于对象数。
\end{itemize}

\subsection{上下文智能降载（按需下探）}
每轮上下文固定携带第 1 层结构化全集、第 2 层卡片检索结果（query 命中或最近书写兜底，上限 K 条，仅摘要）和路径。第 3 层全文与第 4 层资产均不进入上下文。AI 判断需要细节时，通过工具按路径定向读取，读取粒度是单个对象，成本为常数。效果：单轮上下文体积与记忆总量解耦，模型 token 消耗、桥端 IO、内存占用三者同时受控。

\subsection{人类输入信号闭环（任何位置的新信息不被遗漏）}
人类在正文中的每次保存（应用内保存，可与 AI 写入区分）追加一条未确认记录，用追加式日志保存，含路径、指纹、摘要，崩溃安全。AI 每次会话上下文的人类新输入列表由两部分合并，结构标注与未确认正文写入并列呈现、同样必读。AI 引用后经确认接口应答，记录转为已确认，不再重复提醒；同一对象再次被人类修改保存则重新亮起。闭环的保证路径是写入时打标记，读取时必列，确认后不再打扰，与输入载体无关。其状态流转如图~\ref{fig:signal} 所示。

\subsection{数据一致性与崩溃安全（可并入已有公知做法作为配套）}
固定提交顺序为 WAL 落盘并 fsync、命令 reducer 应用、全量校验、临时文件 fsync、原子替换、目录 fsync、返回新版本号。版本冲突不静默合并，基线落后且字段重叠时返回 409，双方数据保留，由人或 AI 显式读取后合并。命令按 commandId 幂等；外部损坏文件移入隔离区，不污染当前视图；崩溃后启动重放未完成 WAL（RPO=0）。

\section{实施方式（关键流程示例）}
\textbf{流程 1，保存即建卡（写入路径）}，如图~\ref{fig:save} 所示：
\begin{enumerate}[leftmargin=2em,itemsep=2pt]
\item 收到正文保存请求（人类或应用），校验，原子写入正文文件；
\item 读回文件，计算标题、摘要、指纹；
\item 合并更新卡片索引并原子写盘；
\item 追加人类输入未确认记录；
\item 向客户端返回新内容与 etag。
\end{enumerate}

\textbf{流程 2，查询下探（读取路径）}，如图~\ref{fig:query} 所示：
\begin{enumerate}[leftmargin=2em,itemsep=2pt]
\item AI 请求上下文与检索词；
\item 读卡片索引（单文件），对摘要场做检索打分，保留可解释理由；
\item 逐对象 lstat 校验指纹，失效者现场重读并刷新该卡；
\item 返回结构化全集、命中摘要（不超过 K 条、含理由）、未确认人类输入和路径；
\item AI 需要细节时使用定向读正文工具或资产路径工具。
\end{enumerate}

\begin{figure}[ht]
\centering
\begin{tikzpicture}[font=\small,
  box/.style={draw=black!70,rounded corners=2pt,minimum width=8.6cm,inner sep=5pt},
  arrow/.style={-,thick}]
  \node[box,fill=blue!12] (s1) {1. 保存请求（人类/应用）$\rightarrow$ 校验 $\rightarrow$ 原子写入正文};
  \node[box,fill=green!12,below=6pt of s1] (s2) {2. 读回文件，计算标题、摘要、指纹};
  \node[box,fill=yellow!22,below=6pt of s2] (s3) {3. 合并更新卡片索引，原子写盘};
  \node[box,fill=orange!18,below=6pt of s3] (s4) {4. 追加人类输入未确认记录};
  \node[box,fill=green!12,below=6pt of s4] (s5) {5. 返回新内容与 etag};
  \draw[arrow] (s1)--(s2); \draw[arrow] (s2)--(s3); \draw[arrow] (s3)--(s4); \draw[arrow] (s4)--(s5);
\end{tikzpicture}
\caption{保存即建卡（写入路径）}\label{fig:save}
\end{figure}

\begin{figure}[ht]
\centering
\begin{tikzpicture}[font=\small,
  box/.style={draw=black!70,rounded corners=2pt,minimum width=8.6cm,inner sep=5pt},
  arrow/.style={-,thick}]
  \node[box,fill=blue!12] (q1) {1. AI 请求上下文与检索词};
  \node[box,fill=green!12,below=6pt of q1] (q2) {2. 读卡片索引（单文件），摘要场检索打分（附理由）};
  \node[box,fill=yellow!22,below=6pt of q2] (q3) {3. 逐对象 lstat 校验指纹，失效者现场重读刷新};
  \node[box,fill=orange!18,below=6pt of q3] (q4) {4. 返回全集+命中摘要($\le$K 条)+未确认输入+路径};
  \node[box,fill=green!12,below=6pt of q4] (q5) {5. 需细节时：定向读正文 / 资产路径工具};
  \draw[arrow] (q1)--(q2); \draw[arrow] (q2)--(q3); \draw[arrow] (q3)--(q4); \draw[arrow] (q4)--(q5);
\end{tikzpicture}
\caption{查询下探（读取路径）}\label{fig:query}
\end{figure}

\begin{figure}[ht]
\centering
\begin{tikzpicture}[font=\small,
  box/.style={draw=black!70,rounded corners=2pt,minimum width=7.2cm,inner sep=5pt},
  arrow/.style={-,thick}]
  \node[box,fill=blue!12] (h1) {外部修改文件（绕过系统）};
  \node[box,fill=yellow!22,below=6pt of h1] (h2) {下次查询：卡片指纹 vs 文件 mtime/size};
  \node[box,fill=red!12,below=6pt of h2] (h3) {不一致（失效）};
  \node[box,fill=green!12,below=6pt of h3] (h4) {只重读该文件，刷新单卡};
  \node[box,fill=green!12,right=14pt of h2] (h5) {一致：直接用摘要（零读取）};
  \draw[arrow] (h1)--(h2); \draw[arrow] (h2)--(h3); \draw[arrow] (h3)--(h4);
  \draw[arrow] (h2)--(h5); \draw[arrow] (h4.south) -- ++(0,-6pt) -| (h5.south);
  \node[below=2pt of h5,font=\scriptsize] {索引一致收敛};
\end{tikzpicture}
\caption{双写与绕桥修改的自愈（指纹失效重读单卡）}\label{fig:heal}
\end{figure}

\begin{figure}[ht]
\centering
\begin{tikzpicture}[font=\small,
  box/.style={draw=black!70,rounded corners=2pt,minimum width=7.4cm,inner sep=5pt},
  arrow/.style={-,thick}]
  \node[box,fill=blue!12] (c1) {人类保存正文（写入时打标记）};
  \node[box,fill=yellow!22,below=6pt of c1] (c2) {未确认记录 new};
  \node[box,fill=orange!18,below=6pt of c2] (c3) {AI 上下文列表必列（与标注并列）};
  \node[box,fill=green!12,below=6pt of c3] (c4) {AI 引用并经确认接口应答 ack};
  \node[box,fill=green!18,below=6pt of c4] (c5) {不再重复提醒（acknowledged）};
  \draw[arrow] (c1)--(c2); \draw[arrow] (c2)--(c3); \draw[arrow] (c3)--(c4); \draw[arrow] (c4)--(c5);
  \draw[arrow,red] (c5.west) -- ++(-1.1,0) |- (c1.west) node[midway,left,font=\scriptsize]{再次修改保存，重新亮起};
\end{tikzpicture}
\caption{人类输入信号闭环}\label{fig:signal}
\end{figure}

\section{技术效果}
\begin{center}
\begin{tabular}{>{\bfseries}p{3.6cm}p{8.8cm}}
\toprule
指标 & 效果（实测）\\
\midrule
AI 每轮上下文体积 & 与记忆总量解耦，500 对象 0 全文读取\\
检索成本 & 常数级，1 个索引文件加 N 次 lstat\\
索引一致性 & 任何写路径或绕桥修改后，一次查询内自愈\\
人类输入丢失率 & 结构上消除，双通道未读信号加确认闭环\\
召回可解释性 & 每次召回附确定性理由\\
信息载体覆盖 & 任意格式附件以原始载体保真可达，不限于文本\\
写入安全 & RPO=0、409 冲突保留、commandId 幂等\\
部署约束 & 零外部依赖、本地优先、文件即数据库\\
\bottomrule
\end{tabular}
\end{center}

\section{创新点（与现有技术的区别，按新颖性强度排序）}
\begin{enumerate}[leftmargin=2em,itemsep=3pt]
\item \textbf{对象资料包为记忆单元，路径化下探，跨格式完整保真。} 以资料包目录为单位组织任意格式的附件（图片、Word、PDF、Excel、多媒体），Markdown 为共通入口，路径为纽带。RAG 只做文本召回或强行转文本，会有信息损失；本方案只要读取方具备对应格式的读取能力，原始载体就完整可达。图片就是图片，Word 就是 Word，信息不因检索机制而丢失。
\item \textbf{图概要、卡片摘要、正文、资产四级分层与上下文按需下探。} 既不用 RAG/向量库把内容压平注入上下文，也不用全文内存检索的无索引方案。AI 上下文只含概要、命中摘要与路径，全文与资产必须显式下探，上下文成本与记忆总量解耦。
\item \textbf{指纹驱动的惰性自愈索引。} 用 lstat 廉价指纹做索引新鲜度判定，不需要每次全量重建，也不依赖外部变更通知；双写与绕桥修改统一自愈，一次查询内收敛。
\item \textbf{人类正文输入信号闭环（md 未读信号加 AI 已读确认）。} 把人对系统的重要输入从依赖检索碰命中的偶然行为提升为结构保证，与结构标注并列、同样必读。
\item \textbf{可解释确定性检索。} 在不引入 embedding 或外部模型的前提下实现可控召回成本，符合本地优先与零依赖约束。
\end{enumerate}

\end{document}